Tempromandibular disorders (TMD) are common with patients experiencing limited temporomandibular joing (TMJ) function.  
This limited function often affects several critical components of daily living.  
Symptoms such as joint sounds and mild pain are often found in both patients with TMJ disorders and those with dentofacial deformities.  
Despite this, the effect of skeletal malocclusion classification (Class II and III) and sex on TMJ function is an area of minimal 
understanding.  Additionally, orthognatic surgery can alleviate some or all of these symptoms.  Mandibular kinematic data reflects 
TMJ functional motion and has vast potential for clinical use.  

Patients enrolled in an IRB-approved prospective natural history study on morphological differences across dentofacial deformities at 
the NIH betweeen 2018 and 2022.  Forty patients with Class I, II, and III (asymptomatic=24, symptomatic=16) completed a biomechanical 
ssessment of TMJ function.  The average age was 32.73 years, with a standard deviation of 14.81 years.  
Twenty-five of the subjects were female, while fifteen were male.  TMJ symptoms include popping or clicking sounds, discomfort, 
and mild pain.  3D manidublar kinematics were captured using a custom jaw tracking system during maximum open-close, lateral 
(left and right), and anterior oral tasks.  

Each subject is fitted with several markers on several parts of the jaw: the lower central incisor as well as the left and right 
apex of the mandible.  Data combined with CBCT images resulted in kinematic trajectories of the mandibular ventral incisor (CI), 
and movement of the bilateral condylar apexes (CA).  For each subject and task, numerous replications are run and three dimensional 
kinematic data in cartesian coordinates relative to the starting position are recorded; for each task, the X, Y, and Z coordinates of 
jaw movement are observed.  In other words, we are provided tri-variate functional outcomes \(P_{ijk}^X(t), P_{ijk}^Y(t), P_{ijk}^Z(t)\) 
for subject \(i\), repetition \(j\), and task \(k\).